@@M0@@ the intersection number, which can be calculated as the coefficient of the monomial @@M1@@ in the Snapper polynomial @@M2@@ (X 4.3.1).

\sgaStableTarget{sga6.e13.para.2-2}\textbf{2.2.} From now on, we suppose that @@M3@@ is algebraically closed, that @@M4@@ is projective, endowed with a very ample invertible sheaf @@M5@@, and that @@M6@@ is integral, of dimension @@M7@@.

Let @@M8@@ be an invertible sheaf on @@M9@@. The integer @@M10@@ is called the \uline{degree} of @@M11@@ (relative to @@M12@@) and is denoted by @@M13@@.

In what follows we consider an invertible sheaf @@M14@@ and the Hilbert polynomials: @@M15@@ and @@M16@@.

\sgaStableTarget{sga6.e13.lem.2-3}\textbf{Lemma 2.3.} \emph{We have}:
\begin{enumerate}[label=(\roman*)]
\item \emph{@@M17@@.}
\item \emph{(If @@M18@@), @@M19@@.}
\end{enumerate}

Indeed, we obtain (i) and (ii) at once by taking @@M20@@ in the following polynomial:
@@M21@@

\sgaStableTarget{sga6.e13.lem.2-4}\textbf{Lemma 2.4.} \emph{If @@M22@@ and @@M23@@, then @@M24@@; therefore @@M25@@ and @@M26@@.}

Indeed, since @@M27@@ is integral, there exists a divisor @@M28@@ such that @@M29@@. If @@M30@@, we would have @@M31@@.

\sgaStableTarget{sga6.e13.lem.2-5}\textbf{Lemma 2.5.} \emph{If @@M32@@, then @@M33@@ is a @@M34@@-sheaf (with @@M35@@ zeros).}

